\documentclass[11pt,a4paper]{article}
\usepackage[T1]{fontenc}
\usepackage[utf8]{inputenc}
\usepackage{lmodern,amsmath,amssymb,amsthm,mathtools}
\usepackage[margin=29mm]{geometry}
\usepackage{enumitem,microtype}
\usepackage[colorlinks=true,linkcolor=blue,citecolor=blue,urlcolor=blue]{hyperref}
\setlength{\emergencystretch}{2em}
\setlist[enumerate]{itemsep=3pt,topsep=5pt}
\setlist[itemize]{itemsep=3pt,topsep=5pt}
\newtheorem{theorem}{Theorem}[section]
\newtheorem{lemma}[theorem]{Lemma}
\newtheorem{proposition}[theorem]{Proposition}
\newtheorem{corollary}[theorem]{Corollary}
\theoremstyle{definition}
\newtheorem{definition}[theorem]{Definition}
\newtheorem{example}[theorem]{Example}
\newtheorem{remark}[theorem]{Remark}
\newcommand{\N}{\mathbb N}
\newcommand{\Ftwo}{\mathbb F_2}
\newcommand{\norm}{\operatorname{N}}
\newcommand{\headpart}{\operatorname{lim}}
\newcommand{\ord}{\operatorname{ord}}
\numberwithin{equation}{section}
\title{Controlled Ordinal Descents, Restricted Partitions\\
and Finite Cartier Observers}
\author{}
\date{17 November 2026\\Working manuscript}
\begin{document}
\maketitle
\begin{abstract}
A terminating ordinal descent can retain difficult modular information
in its accumulated event count. We prove that a polynomial block under
congruent growing budgets has the same count modulo a fixed modulus as
a constant-budget restricted-partition problem. For the standard
finite-hierarchy control in the Wilken--Weiermann setting, a genuine
retained-prefix block yields a partition coefficient whose cutoff and
index are fixed by one entry parameter. We give a complete finite
coefficient observer and prove a closed residue-profile representation
after the pole-order transient. Both binary transitions and the scalar
output are explicit. On the linked inputs, a particular primitive-root
residue tuple has exact period twice one less than the binary logarithm
of the first block budget. This period grows without bound. It rules out
a uniformly bounded cycle for the complete rational state, while leaving
the prescribed scalar observation and unchanged-control canonical
comparison open. All predecessor, budget-transfer, observer, and period
arguments needed for these results are included.
\end{abstract}
\begingroup
\small
\tableofcontents
\endgroup
\clearpage

\section{Counting beyond termination}

Norm-controlled ordinal assignments can turn termination proofs into
numerical complexity bounds. In the assignment for G\"odel's T studied
by Wilken and Weiermann \cite{ww}, a collapse value can be evaluated by a
strictly decreasing ordinal descent. Its norm budget grows rapidly, and
the number of replacement events contributes to the returned natural
number. Termination of each evaluation does not decide a quantified
question about those event counts. Rapid growth by itself does not
determine their small residues.

This article studies that arithmetic for the unchanged control
\[
 F_0(n)=2^n,\qquad F_{j+1}(n)=F_j^{\,n+1}(n),\qquad
 G(n)=F_5(n+100).
\]
Superscripts on functions denote iteration. The general transfer theorem
needs only congruence hypotheses on the budgets; the concrete
retained-prefix family and its period calculation use this particular
$G$.

There are three different finite objects in the argument. For a supplied
cutoff $d$, a restricted-partition generating function has a complete
finite numerator observer. After enough binary selections, it has an
equivalent complete vector of simple residues in a finite splitting
field. A selected family of primitive roots provides a projection of
that vector, sufficient to prove a period lower bound but insufficient
to recover the output coefficient. Keeping these objects distinct is
essential.

\begin{theorem}[Main results]\label{thm:main}
The constructions below have the following properties.
\begin{enumerate}
\item Let $d,m\ge1$, let $p_d(n)$ count partitions of $n$ with parts at most $d$, and
put $Q_d=d!m^{d-1}$. A complete deletion of one $\omega^d$ suffix
with retained mass $c\ge0$ and initial total norm at least $c+d+1$
has event count congruent to $p_d(S-c)$ modulo
$m$, whenever every operational budget is at least the current norm
and is congruent to $S\ge c$ modulo $Q_d$.
\item For each fixed $d$, the binary coefficients of
$\prod_{i=1}^d(1+z^i)^{-1}$ are computed exactly by two explicit
linear maps on a space of dimension $d(d+1)/2$ over $\Ftwo$.
Distinct numerator states have distinguishing continuations of length
at most $\lceil\log_2(d(d+1)/2)\rceil$.
\item After an explicit finite transient, the complete residue profile
is closed under both binary selectors. At a pole $\zeta$ its two
updates are $b_\zeta\mapsto b_\zeta^{1/2}$ and
$b_\zeta\mapsto(\zeta b_\zeta)^{1/2}$; the scalar output is
$\sum_\zeta b_\zeta/\zeta$.
\item For every $v\ge4$, the descent from
$\omega^{\omega+1}+(v-4)$ has an explicit first
$\omega^\omega$-block parity. On its linked partition parameters,
a primitive-root residue tuple has exact period
$2(\log_2 B-1)$, where $B=G(G(v))$.
Every eventual period of the complete differential Cartier state
is divisible by this number.
\end{enumerate}
\end{theorem}

The first item is a statement about complete segment counts. The last
item is a statement about complete states and a particular residue
projection. Neither states that the selected scalar has the same period.
The displayed ordinal input family is actual syntax for the collapse
defined below. An interpretation of this family inside a particular
larger typed-term assignment is a separate source question.

\section{Norms, maximum predecessors and exact counts}
\label{sec:descent}

All ordinals lie below $\varepsilon_0$ and are represented by hereditary
Cantor normal forms. Addition and multiplication by finite coefficients
have their ordinary ordinal meanings. For strictly decreasing exponents
and positive finite coefficients define
\[
 \norm(0)=0,\qquad
 \norm\!\left(\sum_{i<r}\omega^{\alpha_i}c_i\right)
       =\sum_{i<r}c_i(1+\norm(\alpha_i)).
\]
Thus $\norm(n)=n$, $\norm(\omega^d)=d+1$ for finite $d$,
$\norm(\omega^\omega)=3$, and
$\norm(\omega^{\omega+1})=4$. This is the additive norm used
in the relevant norm-controlled hierarchy framework
\cite{bcw,ww}.

\begin{lemma}[Effective norm balls]\label{lem:norm-balls}
Each norm ball is finite and effectively enumerable.
Cantor comparison and the norm are computable.
\end{lemma}
\begin{proof}
For budget zero there is only zero. In a nonzero ordinal of norm at
most $b$, every exponent has norm at most $b-1$. Inductively enumerate
those exponents, order them by recursive Cantor comparison, and enumerate
the finite coefficient lists whose total charge is at most $b$.
This produces every candidate and terminates. The same recursion
computes the norm and compares the resulting forms.
\end{proof}

Fix for now a computable strictly increasing $f:\N\to\N$ with
$f(n)>n$. Define
\[
 \psi_f(\alpha)=\max\bigl(\{0\}\cup
   \{\psi_f(\beta)+1:\beta<\alpha,\ 
                    \norm(\beta)\le f(\norm(\alpha))\}\bigr).
 \tag{2.1}
\]
Ordinal recursion is legitimate because the maximum is finite.
Its defining controlled implication is
\[
 \beta<\alpha,\quad \norm(\beta)\le f(\norm(\alpha))
 \quad\Longrightarrow\quad
 \psi_f(\beta)+1\le\psi_f(\alpha).
 \tag{2.2}
\]
Induction gives $\psi_f(n)=n$ for finite $n$.
No unrestricted monotonicity of $\psi_f$ is assumed.

\subsection{The exact maximum predecessor}
For $\alpha>0$ and $b\in\N$ put
\[
 M(\alpha,b)=\max\{\beta<\alpha:\norm(\beta)\le b\},
 \qquad L(e,b)=M(\omega^e,b).
\]
The candidate set contains zero. Put $M(0,b)=0$ as a convention.

\begin{theorem}[Greedy predecessor]\label{thm:greedy}
The following simultaneous recursion computes $M,L$ exactly.
Write $\alpha=\omega^e c+r$ with $c>0$, $r<\omega^e$, and
$w=1+\norm(e)$.
\begin{enumerate}[label=(\roman*)]
\item $M(\alpha,0)=0$ and $L(0,b)=L(e,0)=0$.
\item If $r>0$ and $cw\le b$, then
\[
 M(\alpha,b)=\omega^e c+M(r,b-cw).
\]
\item Otherwise, put
\[
 a=\begin{cases}
 \min(c-1,\lfloor b/w\rfloor),&r=0,\\
 \min(c,\lfloor b/w\rfloor),&r>0.
 \end{cases}
\]
Then $M(\alpha,b)=\omega^e a+L(e,b-aw)$.
\item If $e,b>0$, put $\eta=M(e,b-1)$,
$v=1+\norm(\eta)$ and $(q,s)=\operatorname{divmod}(b,v)$.
Then $L(e,b)=\omega^\eta q+L(\eta,s)$.
\end{enumerate}
Zero blocks are omitted.
\end{theorem}
\begin{proof}
When the whole leading block and a nonzero tail bound are affordable,
retaining that block and lowering the tail dominates every predecessor
with a smaller leading block. This proves (ii). Otherwise (iii) chooses
the largest affordable coefficient compatible with being below
$\alpha$; its remainder is strictly below $\omega^e$.

A nonzero ordinal below $\omega^e$ with budget $b$ has a leading
exponent $\eta'<e$ of norm at most $b-1$. Its largest possible
exponent is therefore $M(e,b-1)$. Cantor order next maximizes its
coefficient and then its lower tail, proving (iv).

The lexicographic measure (budget, phase), with $M$ above $L$,
proves termination. Clause (ii) lowers the budget. Clause (iii)
either lowers it or changes phase from $M$ to $L$. In (iv) the
first call uses $b-1$; also $v\le b$, so $q\ge1$ and $s<b$
for the second call. The value returned in each terminating case
is the greatest candidate just characterized.
\end{proof}

\begin{lemma}[Saturation]\label{lem:saturation}
For every nonzero limit $\lambda$ and every budget $b$,
$\norm(M(\lambda,b))=b$. The maximum is strictly increasing
with $b$.
\end{lemma}
\begin{proof}
If $\eta=M(\lambda,b)$ left unused budget
$a=b-\norm(\eta)>0$, the larger ordinal $\eta+a$ would still
be below the limit $\lambda$ and would have norm $b$, a
contradiction. If $b<c$, the ordinal $M(\lambda,b)+1$
is below $\lambda$ and has norm at most $c$, proving strict
increase.
\end{proof}

\begin{theorem}[Single-predecessor recurrence]\label{thm:predecessor}
For every nonzero limit $\lambda$ and $n\in\N$,
\[
 \psi_f(\lambda+n)
 =\psi_f\bigl(M(\lambda,f(\norm(\lambda)+n))\bigr)+1.
\]
\end{theorem}
\begin{proof}
Induct on $n$. Put $b=f(\norm(\lambda)+n)$ and
$\eta=M(\lambda,b)$. Saturation shows that $\eta$ is an admitted
predecessor of norm $b$, so it gives the lower inequality.
Every admitted $\beta<\lambda$ satisfies $\beta\le\eta$.
If $\beta<\eta$, then
$\norm(\beta)\le b=\norm(\eta)<f(\norm(\eta))$,
and controlled strictness bounds its contribution by
$\psi_f(\eta)$.

The other possible predecessors are $\lambda+j$ with $j<n$.
Put $b_j=f(\norm(\lambda)+j)$ and
$\eta_j=M(\lambda,b_j)$. Strict increase of $f$ and saturation
give $\eta_j<\eta$ and $\norm(\eta_j)=b_j<b$.
The induction hypothesis and controlled strictness yield
$\psi_f(\lambda+j)=\psi_f(\eta_j)+1\le\psi_f(\eta)$.
Thus each contribution to the defining maximum is at most
$\psi_f(\eta)+1$, proving the reverse inequality.
\end{proof}

\subsection{The complete event count}
For a nonfinite ordinal $\alpha$, write
$\alpha=\headpart(\alpha)+n$, where the first summand is a
nonzero limit and $n$ is finite. At each event replace
\[
 \alpha\longmapsto M(\headpart(\alpha),f(\norm(\alpha))).
 \tag{2.3}
\]
In particular the predecessor bound is the limit part, not the
original ordinal with its finite tail.

For a fixed limit $\lambda$ and total input norm
$t\ge\norm(\lambda)$, start at $\lambda+t-\norm(\lambda)$.
Let $\tau_\lambda(t)$ count events until the first finite
ordinal $Z_\lambda(t)$, and put
$J_\lambda(t)=\psi_f(\lambda+t-\norm(\lambda))$.

\begin{proposition}[Complete-count identity]\label{prop:complete}
The descent terminates and
\[
 Z_\lambda(t)=f^{\,\tau_\lambda(t)}(t),\qquad
 J_\lambda(t)=Z_\lambda(t)+\tau_\lambda(t).
\]
\end{proposition}
\begin{proof}
Every event strictly decreases the ordinal. Well-foundedness below
$\varepsilon_0$ therefore proves termination. Saturation gives the
successive norms $t,f(t),f^2(t),\ldots$; the finite endpoint equals
its own norm. Theorem~\ref{thm:predecessor} adds one for each
event, and finite evaluation returns $Z_\lambda(t)$.
\end{proof}

There is no further event count through finite successors below the
endpoint. A count stopped at an earlier nonfinite boundary is a
different segment count. The block calculations below always specify
which prefix remains and when their suffix first becomes finite.

\section{The modular information supplied by the control}
\label{sec:control}

For the remaining concrete applications fix $f=G$, where
\begin{equation}\label{eq:control}
 F_0(n)=2^n,\qquad F_{j+1}(n)=F_j^{\,n+1}(n),\qquad
 G(n)=F_5(n+100).
\end{equation}
The budget-transfer result in the next section remains a theorem for
arbitrary schedules satisfying its own premises.

\begin{lemma}[Hierarchy growth]\label{lem:growth}
Every $F_j$ is strictly increasing and $F_j(n)>n$. For $n\ge1$,
$F_j(n)\ge2^n$. The values of $G$ are powers of two, and
\begin{equation}\label{eq:growth}
 G(n)\ge2^{n+100},\qquad G(n+1)\ge2^{G(n)+1}.
\end{equation}
\end{lemma}
\begin{proof}
The assertions of strict increase and inflation follow simultaneously
by induction on $j$: increasing the seed increases every common
iterate, and the larger iteration count supplies another inflationary
step. The bound by $2^n$ follows from the first lower-level call.
Every final call ultimately ends with $F_0$, giving a power of two.

Put $x=n+100$ and $A=F_5(x)=F_4^{x+1}(x)$. Then
$F_4^{x+1}(x+1)\ge A+1$, and
\[
 G(n+1)=F_4\bigl(F_4^{x+1}(x+1)\bigr)
       \ge F_4(A+1)\ge2^{A+1}.
\]
The other bound is the finite-level domination already proved.
\end{proof}

We state the arithmetic input in a form that can be applied when a
modulus itself depends on a source parameter.

\begin{theorem}[Uniform residue envelopes]\label{thm:envelope}
There is an effective function $m\mapsto r_m\in\{0,\ldots,m-1\}$ such
that:
\begin{enumerate}
\item $f(x)\equiv r_m\pmod m$ for $x\ge m$;
\item for $s\ge1$, $j\ge2$ and $m\le2^{2^s}$,
      $F_j(s)\equiv r_m\pmod m$;
\item if $B=f(u)$, then
      $f(u+1)\equiv r_m\pmod m$ for every $m\le2^{2^B}$.
\end{enumerate}
\end{theorem}
\begin{proof}
We recall the construction to make the envelopes precise. Put $r_1=0$.
For $m>1$, write $m=2^a b$, $b$ odd, and set
$q=\varphi(b)$ if $b>1$, otherwise $q=1$. Then $q<m$.
Define $r_m$ by
$r_m\equiv0\pmod {2^a}$ and
$r_m\equiv2^{r_q}\pmod b$. This is a terminating recursion using
Euler's theorem and the Chinese remainder theorem.

Let $E_0(n)=n$, $E_{h+1}(n)=2^{E_h(n)}$.
Induction on $m$ proves that $E_h(n)\equiv r_m\pmod m$ whenever
$h,n\ge m$. The exponent is large enough for the $2$-part; its residue
modulo $q$ is supplied by the induction hypothesis, since $h-1\ge q$.
As $F_1(n)=E_{n+1}(n)$, every $F_j(n)$ with $j\ge1$ has residue
$r_m$ for $n\ge m$: in a higher-level iterate every argument is at
least $n$. This proves the first assertion.

For the stronger bound, the final $F_1$ call in
$F_2(s)=F_1^{s+1}(s)$ has argument at least
$F_1(s)\ge2^{2^s}$ when $s\ge1$. Its tower seed and height are both
at least $m$, proving the second assertion for $j=2$.
The induction to higher $j$ uses its last lower-level call, whose
argument is at least $s$.

Finally put $n=u+100$. In
$F_5(n+1)=F_4^{n+2}(n+1)$ the argument of the final $F_4$ is
$v=F_4^{n+1}(n+1)>F_4^{n+1}(n)=B$. The second assertion at $v$
therefore gives the third.
\end{proof}

The theorem does not say that an arbitrary computable sequence of
moduli may be stabilized at an arbitrary input. Each use must pay the
displayed bound. In particular, a lower bound on $f(x)$ by itself is
not the first assertion's lower bound on its \emph{argument}.

The residues are compatible under divisibility: if $a\mid b$, evaluate
$G(x)$ at any $x\ge b$ to obtain $r_b\equiv r_a\pmod a$.
For the concrete entry calculation below, $r_2=0$ and
$r_3=2^{r_2}\bmod3=1$.


\section{Constant budgets and restricted partitions}

For $d\ge1$, define
\[
 p_d(n)=\#\{(a_1,\ldots,a_d)\in\N^d:
                         a_1+2a_2+\cdots+da_d=n\}.
\]
The independent choices of multiplicities give
\begin{equation}\label{eq:partition-series}
 \sum_{n\ge0}p_d(n)z^n=\prod_{i=1}^d(1-z^i)^{-1}.
\end{equation}
The relation to ordinals preserves every part, including the finite tail.

\begin{lemma}[A fixed-budget count]\label{lem:fixed}
Fix a nonnegative retained mass $c$ and take $S\ge c+d+1$.
Start the suffix at $\omega^d$ with any finite tail of total norm
(including $c$) at most $S$. At every replacement spend the constant
suffix budget $S-c$ and stop when the suffix is finite. The number of
replacements is $p_d(S-c)$. When $c=\norm(\pi)$ for a retained
CNF prefix $\pi$ of least exponent at least $d$, this is precisely
the corresponding segment starting at $\pi+\omega^d$ plus that tail.
\end{lemma}
\begin{proof}
The ordinals of suffix norm exactly $S-c$ below $\omega^d$ are
precisely
\[
 \omega^{d-1}a_d+\cdots+\omega a_2+a_1,\qquad
             \sum_{i=1}^d i a_i=S-c.
 \tag{*}
\]
Their number is $p_d(S-c)$. The first predecessor is the largest
such ordinal. In any later state of (*), the nonfinite coefficients
determine $a_1$ uniquely. Thus no second state of (*) has the same
nonfinite part. Removing the finite tail and taking the largest
predecessor of suffix norm $S-c$ selects exactly the next state of
(*) in decreasing ordinal order. All states are visited, ending at
the finite suffix $S-c$. There is one replacement into each state
and none after the last, proving the count. For $d=1$ there is one
state and one replacement.
\end{proof}

\begin{lemma}[A sufficient partition period]\label{lem:period}
For every $d,m\ge1$, put $Q_d=d!m^{d-1}$. Then
\[
             p_d(n+Q_d)\equiv p_d(n)\pmod m\quad(n\ge0).
\]
\end{lemma}
\begin{proof}
For $d=1$, $p_1(n)=1$. For $d>1$, splitting by the number
of parts of size $d$ gives
$p_d(n)=\sum_{j=0}^{\lfloor n/d\rfloor}p_{d-1}(n-dj)$.
Here $Q_d/d=mQ_{d-1}$. In the sum for $n+Q_d$, shifting the
index by $mQ_{d-1}$ recovers the old sum. The additional
$mQ_{d-1}$ summands form $m$ groups of $Q_{d-1}$ consecutive
indices in the periodic sequence
$j\mapsto p_{d-1}(n+Q_d-dj)\pmod m$.
Its period divides $Q_{d-1}$ by induction, so the extra sum is zero
modulo $m$.
\end{proof}

The period is deliberately a sufficient bound, not an optimal one.
For $d=2$ the count is $\lfloor n/2\rfloor+1$, and $2m$
is the period supplied by the lemma.

\section{A transfer theorem for congruent growing budgets}

For the numerical statement, a retained mass $c\ge0$ means that a
suffix of norm $n$ has total norm $c+n$, and a total budget $b$ leaves
suffix budget $b-c$. It need not be the norm of an ordinal prefix.
If it is realized by an actual retained CNF prefix, the numerical
process agrees with that source segment. A budget schedule may increase
with the current total norm; it must always afford that norm.

\begin{theorem}[Congruent-budget transfer]\label{thm:transfer}
Consider deletion of one suffix block $\omega^d$, with retained
numerical mass $c\ge0$ and initial total norm at least $c+d+1$,
under budgets at least the current total norm. Suppose
every budget used in this segment is congruent to one integer $S$
modulo $Q_d=d!m^{d-1}$, where $S\ge c$. Then its complete segment
count is congruent to $p_d(S-c)$ modulo $m$.
\end{theorem}
\begin{proof}
Induct on $d$. A degree-one block takes one replacement, as does
$p_1(S-c)$. For the induction step let $b$ be the first actual
total budget. The suffix of its greatest predecessor is the greedy
Cantor polynomial of norm $b-c$ and degree less than $d$.
Write its coefficients as $a_{d-1},\ldots,a_0$, with the final
finite coefficient included.

Thereafter each of its nonfinite monomial units is deleted by a
lower-degree block. For a degree-$i$ unit the retained prefix norm
has the form
\[
 c_i=c+\sum_{j>i}(j+1)a_j+(i+1)a,\qquad 0\le a<a_i.
 \tag{**}
\]
The order of those units is fixed by the Cantor order; growing
budgets affect their lower suffixes and durations, not the number
of these fixed higher monomial units. Since $Q_i\mid Q_d$, every
budget in a child segment has the residue of $b$ modulo $Q_i$.
Induction makes its count congruent to $p_i(b-c_i)$.
These arguments are nonnegative: each $c_i$ consists of $c$ plus
the norm of a prefix of the first suffix polynomial and is at most $b$.

Consequently the whole count, including its initial replacement,
has the same residue as the constant-$b$ process with exactly that
initial polynomial and those prefixes (**). By
Lemma~\ref{lem:fixed} this is $p_d(b-c)$.
Since $b\equiv S\pmod {Q_d}$, Lemma~\ref{lem:period} gives
$p_d(b-c)\equiv p_d(S-c)\pmod m$. If $b-c<S-c$,
apply the period lemma in the reverse direction by adding the
appropriate nonnegative multiple; no negative partition index is used.
\end{proof}

The proof covers the whole count, not just its leading contribution.
It does not replace the increasing operational budgets by a constant
control: the latter is only a device for computing one residue.

\begin{corollary}[A fixed polynomial block under the standard control]
\label{cor:fixed-polynomial-residue}
Fix $c,d,m$ and put $Q_d=d!m^{d-1}$. Under the control $G$,
the complete count $E_c^d(t)$ of one $\omega^d$ block with retained
mass $c$ and total input norm $t$ satisfies
\[
 E_c^d(t)\equiv
 p_d\bigl((r_{Q_d}-c)\bmod Q_d\bigr)\pmod m
 \quad\text{if }t\ge\max(c+d+1,Q_d).
\]
The finitely many smaller valid inputs can be evaluated exactly.
\end{corollary}
\begin{proof}
All later norms are at least $t$, so every budget has residue
$r_{Q_d}$ by Theorem~\ref{thm:envelope}. Apply the transfer theorem
with $S=c+((r_{Q_d}-c)\bmod Q_d)\ge c$.
Exact evaluation handles the remaining finite set.
\end{proof}


\section{A retained prefix and its linked partition coefficient}
\label{sec:boundary}

\begin{proposition}[An actual $\omega^\omega$ block]\label{prop:omega-block}
Let $\Omega=\omega^\omega$. At a real suffix state $\pi+\Omega+n$, with $\pi$ a retained CNF
prefix (zero or of least exponent at least $\omega$), and total norm $u$,
assume $c=\norm(\pi)$, $c+3\le u$, and
$B=f(u)\ge\max(16,m)$. Put $C=f(u+1)$ and $d=B-c-1$.
If $H_c(u)$ counts the complete deletion of this $\Omega$ copy,
then
\begin{equation}\label{eq:omega-block}
 H_c(u)\equiv1+p_{B-c-1}(C-c)\pmod m.
\end{equation}
Its output norm is still exactly $f^{H_c(u)}(u)$.
\end{proposition}
\begin{proof}
The first predecessor is $\pi+\omega^d$, of total norm $B$.
Indeed the greatest exponent of norm at most $B-c-1$ is the
finite exponent $d$ in this particular bound below $\omega$;
this is the predecessor of $\omega^\omega$ under the suffix budget.
All subsequent norms $x$ are at least $B$. Their budgets
$G(x)=F_5(x+100)$ have residue $r_{Q_d}$ by the second part of
Theorem~\ref{thm:envelope}, because
\[
 Q_d=d!m^{d-1}\le B^{2B}\le2^{2^B}.
\]
The last inequality holds for $B\ge16$; it follows from
$2B\log_2 B\le2^B$, first at $16$ and then by the increasing
ratio of the right side to the left side.
The adjacent value $C=f(u+1)$ has the same residue modulo $Q_d$
by the third part of Theorem~\ref{thm:envelope}.
Apply Theorem~\ref{thm:transfer} to the remaining $\omega^d$ block
with $S=C$, then include its preceding $\Omega$ replacement.
Norm saturation proves the final assertion independently of the
modular calculation.
\end{proof}

\subsection{The first moving-prefix boundary}

Write $\Theta=\omega^{\omega+1}$, whose norm is four.
For each $v\ge4$, begin an actual descent at $\Theta+(v-4)$,
of total norm $v$. The same calculation applies whenever a larger
descent genuinely reaches this limit part at that norm; it does not
posit independent entry values inside a larger source.
Put $u=G(v)$. The residue recursion gives $u\equiv1\pmod3$.

\begin{lemma}[The retained prefix at first entry]\label{lem:entry}
The predecessor of this $\Theta$ state is
\[
 \Omega q+1,\qquad q=(u-1)/3.
\]
The first $\Omega$ copy deleted in this polynomial has retained
prefix norm $e=u-4$ and total input norm $u$.
\end{lemma}
\begin{proof}
A copy of $\Omega$ has norm three. The maximal affordable exponent
below $\omega+1$ is $\omega$, so the greedy predecessor takes as
many copies of $\Omega$ as possible. Since $u\equiv1\pmod3$,
exactly one norm unit remains, whose largest affordable ordinal
is the finite ordinal one. Before deleting the last displayed $\Omega$ in the
ordinary sum $\Omega(q-1)+(\Omega+1)$, the retained prefix is
$\Omega(q-1)$ and has norm $3(q-1)=u-4$.
\end{proof}

With $B=f(u)$, $C=f(u+1)$, $e=u-4$ and $d=B-u+3$, the
first block parity is consequently
\begin{equation}\label{eq:delta}
 \delta_4(v):=H_{u-4}(u)\bmod2
       =1+p_d(C-e)\pmod2.
\end{equation}
The subscript records the norm four of its suffix $\Omega+1$.
The initial $\Theta$ replacement lies outside $H_{u-4}(u)$;
the displayed $1$ instead counts the $\Omega$ replacement.

The value $C=G(u+1)$ is a modular surrogate. It is not the next
operational norm $G(B)$. The transfer theorem permits it because
its residue modulo the required partition period agrees with every
later budget. All four parameters $B,C,e,d$ are determined by the
single entry norm $v$.

This quantity is inside a larger correlated computation. Before its
remaining limit part has become exactly $\Omega$, the $\Theta$ segment
deletes $q-1$ copies with prefix masses $3(q-1),3(q-2),\ldots,3$. Their inputs
are successive points on the \emph{same} $f$ orbit. Observing the
sum of all these counts does not automatically define the graph of
the first summand. Thus \eqref{eq:delta} is a concrete benchmark for
a proposed arithmetic method, not a proved necessary subproblem of
every possible comparison procedure. The family
$\Theta+(v-4)$ proves genuine ordinal-descent realization for every
$v\ge4$; no stronger typed-source grammar realization is used here.


\section{The complete binary coefficient observer}
\label{sec:observer}

Work over $\Ftwo$ and set
\[
 A_d(z)=\prod_{i=1}^d(1+z^i),\qquad D=d(d+1)/2.
\]
For $b\in\{0,1\}$ let
$\Lambda_b(\sum a_nz^n)=\sum a_{2n+b}z^n$.
The state space of proper numerators is
$V_d=\{U\in\Ftwo[z]:\deg U<D\}$.

\begin{lemma}[Exact state update]\label{lem:matrix}
For $U\in V_d$,
\[
 \Lambda_b(U/A_d)=\frac{\Lambda_b(UA_d)}{A_d}.
\]
The updated numerator is again in $V_d$. In the coefficient basis its
matrix is
$T_b[i,j]=[z^{2i+b-j}]A_d$, for $0\le i,j<D$.
\end{lemma}
\begin{proof}
Since $A_d(z)^2=A_d(z^2)$ in characteristic two,
$U/A_d=(UA_d)/A_d(z^2)$. Coefficient selection commutes with a
factor in $z^2$, giving the identity. The numerator degree is at
most $\lfloor(2D-1)/2\rfloor=D-1$. Applying selection to each
monomial basis vector gives the matrix entries.
\end{proof}

\begin{theorem}[Complete selected coefficient]\label{thm:observer}
If $n=\sum_{i<L}b_i2^i$, with bits read least significant first,
then
\[
 p_d(n)\bmod2=e_0^{T}T_{b_{L-1}}\cdots T_{b_0}e_0.
\]
\end{theorem}
\begin{proof}
The initial numerator is one, represented by $e_0$. Successive
selections with $b_0,b_1,\ldots$ discard these bits of the
coefficient index. After all bits, the desired coefficient is the
constant term. As $A_d(0)=1$, this is the constant numerator
coordinate selected by $e_0^T$.
\end{proof}

This formula retains all factors and all cancellations. Boolean
reachability in its transition graph would be a different observation:
two paths contributing to one coefficient cancel in $\Ftwo$.

\begin{lemma}[All-continuation observability]\label{lem:observable}
For any field $K$, polynomial $A$ with $A(0)=1$ and $\deg A=D$,
the first $D$ coefficients of $U/A$ determine every $U$ of degree
less than $D$. For the binary observer above, distinct numerator
states are distinguished by some continuation of length at most
$\lceil\log_2 D\rceil$.
\end{lemma}
\begin{proof}
Multiplication by $1/A=1+a_1z+\cdots$ sends the coefficient
vector of $U$ to the first $D$ output coefficients by a lower
triangular matrix with diagonal one. It is invertible. Equivalently,
the first nonzero coefficient of a difference $U-V$ remains its
first nonzero coefficient after division by $A$. Its index $k<D$
has at most $\lceil\log_2D\rceil$ bits; those bits form a
distinguishing continuation. Index zero permits the empty word.
\end{proof}

This is the linear observability viewpoint familiar from weighted
automata \cite[Sections~1--3]{kiefer}. It does not say that every distinguishing
continuation is available from the genuine source parameters.
A quotient preserving only the one prescribed source word remains a
different question.

\begin{remark}[Linear algebra in characteristic two]
The proof uses the invertible triangular coefficient map directly.
A Gram-matrix argument based on positive definiteness would be invalid
over $\Ftwo$: the nonzero column $v=(1,1)^T$ has rank one but
$v^Tv=0$. Likewise, a large ambient observer dimension does not by
itself prove that all of its states are reachable from numerator one.
\end{remark}


\section{Differential Cartier and the surviving residues}
\label{sec:cartier}

It is useful to normalize the same selected coefficient as a differential.
Over a perfect field $K$ of characteristic two, every rational function
has a unique expression $a(z)^2+zb(z)^2$: the elements $1,z$ form a
basis of $K(z)$ over $K(z^2)$, and perfectness supplies the square
roots of its coefficients. Define
$\mathcal C((a^2+zb^2)\,dz)=b\,dz$.
This is the differential Cartier operator. Multiplication by a square
pulls out its square root.

\begin{lemma}[Local rule]\label{lem:local}
For a fixed pole $\zeta$ and $b\in K$,
\[
 \mathcal C\!\left(\frac{b\,dz}{(z-\zeta)^{2j}}\right)=0,\qquad
 \mathcal C\!\left(\frac{b\,dz}{(z-\zeta)^{2j+1}}\right)
       =\frac{b^{1/2}\,dz}{(z-\zeta)^{j+1}}.
\]
The same definition gives
$\mathcal C(\sum a_nz^n\,dz)=\sum a_{2n+1}^{1/2}z^n\,dz$.
\end{lemma}
\begin{proof}
An even power of $h=z-\zeta$ is a square and contributes zero.
For odd powers, write $h^{2k+1}=(h^k)^2(z-\zeta)$.
The constant differential $\zeta\,dz$ has Cartier image zero,
whereas $\mathcal C(z\,dz)=dz$. Pulling out the square proves
the rule for all integer $k$, including negative $k$.
The coefficient rule follows by the same even/odd decomposition.
\end{proof}

For $e\ge1$ and $2^\ell\ge e$,
\begin{equation}\label{eq:selection}
 [z^0\,dz]\,\mathcal C^\ell\!
             \left(\frac{z^{e-1}}{A_d(z)}\,dz\right)
                   =p_d(2^\ell-e)\pmod2.
\end{equation}
Indeed $\ell$ odd selections choose original exponent $2^\ell-1$,
which in $z^{e-1}/A_d$ has coefficient $p_d(2^\ell-e)$.

\begin{lemma}[Removal of every higher pole]\label{lem:poles}
Put $\mu=\sum_{i=1}^d2^{\nu_2(i)}$ and
$a=\lceil\log_2\mu\rceil$. After $a$ Cartier applications to
a proper differential with denominator $A_d$, all surviving finite
poles are simple. Its simple residues subsequently undergo inverse
Frobenius at their fixed poles. No polynomial part is created.
\end{lemma}
\begin{proof}
The multiplicity at a root of odd order $q$ is
$\sum_{q\mid i,\ i\le d}2^{\nu_2(i)}\le\mu$.
By Lemma~\ref{lem:local}, a term of order $k>1$ survives $a$
applications only if $k-1$ is a positive multiple of $2^a$.
This is impossible for $k\le\mu\le2^a$. A simple residue is
replaced by its square root. The proper partial-fraction expansion
has no polynomial part, and every term in its Cartier image is again
proper, proving the final assertion.
\end{proof}

Here $\nu_2(i)$ denotes the exponent of two in the positive integer
$i$. For rational inputs over $\Ftwo$, the simple-pole tail is
eventually periodic: all its finitely many poles and residues lie
in one finite extension of $\Ftwo$, and inverse Frobenius is a
permutation of that finite field. This proves a finite cycle for
each supplied cutoff, with no cutoff-independent bound.

\subsection{Closure under both binary selectors}

The preceding lemma concerns repeated $\mathcal C$. To obtain a
complete observer for arbitrary remaining binary words, one must
also account for $\Lambda_0$.
For a rational function $R$ with coefficients in $\Ftwo$, define
\[
 \mathcal L_1(R\,dz)=\mathcal C(R\,dz),\qquad
 \mathcal L_0(R\,dz)=\mathcal C(zR\,dz).
 \tag{L}
\]
Their power-series expansions are respectively
$\Lambda_1(R)\,dz$ and $\Lambda_0(R)\,dz$, since square roots
fix the coefficients in $\Ftwo$.

\begin{lemma}[A uniform pole-order transient]\label{lem:binary-transient}
For a proper rational differential with denominator $A_d$,
each $\mathcal L_b$ preserves properness and introduces no
new finite pole. A pole of order $k$ acquires order at most
$\lceil k/2\rceil$. Consequently every binary word of length
$a=\lceil\log_2\mu\rceil$ reaches a proper differential
with only simple poles.
\end{lemma}
\begin{proof}
The assertion for $\mathcal L_1$ is the local Cartier rule.
At a pole $\zeta$, put $h=z-\zeta$. For $\mathcal L_0$,
write
\[
 \frac{bz}{h^k}\,dz
 =\frac{b}{h^{k-1}}\,dz+\frac{\zeta b}{h^k}\,dz.
\]
For $k=2j$ only the first summand survives Cartier, with
order $j$. For $k=2j+1$ only the second survives, with
order $j+1$. When $k=1$, the first summand is a constant
differential and is killed. Multiplication of a proper
rational function by $z$ adds at most a constant
polynomial part, which is likewise killed. Thus the
result is proper with the claimed poles and orders.

Iterating gives order at most $\lceil k/2^a\rceil$.
Every initial pole order is at most $\mu$, and $2^a\ge\mu$.
\end{proof}

Let $\mathcal R_d$ be the set of distinct roots of $A_d$ in
a finite splitting field $K$. These roots are nonzero and
the set is closed under Frobenius $\zeta\mapsto\zeta^2$.
Both the field and the root set can be computed by finite
polynomial factorization and adjoining roots. Define
\[
 \mathcal P_d=
 \left\{(b_\zeta)_{\zeta\in\mathcal R_d}\in K^{\mathcal R_d}:
              b_{\zeta^2}=b_\zeta^2\right\}.
 \tag{P}
\]

\begin{theorem}[A complete residue-profile observer]
\label{thm:residue-observer}
Proper rational differentials over $\Ftwo$ with simple poles
in $\mathcal R_d$ correspond bijectively to $\mathcal P_d$ by
\[
 R(z)\,dz=\sum_{\zeta\in\mathcal R_d}
                 \frac{b_\zeta}{z-\zeta}\,dz.
 \tag{PF}
\]
The profile space is finite and is closed under both binary
transitions. At each fixed pole their updates and the output are
\[
 \mathcal L_1:\ b_\zeta\longmapsto b_\zeta^{1/2},
 \qquad
 \mathcal L_0:\ b_\zeta\longmapsto(\zeta b_\zeta)^{1/2},
 \qquad
 \operatorname{out}(b)=\sum_{\zeta\in\mathcal R_d}
                             \frac{b_\zeta}{\zeta}\in\Ftwo.
 \tag{U}
\]
After the finite transient of
Lemma~\ref{lem:binary-transient}, this is an exact observer
for every remaining binary continuation.
\end{theorem}
\begin{proof}
A proper rational function with these simple poles has the
unique partial-fraction expansion (PF). Applying the field
Frobenius to its coefficients fixes the function precisely
when $b_{\zeta^2}=b_\zeta^2$. Conversely a compatible tuple
makes (PF) invariant under Frobenius, so its coefficients lie
in $\Ftwo$. There is no additional polynomial part.
This proves the bijection.

The local rule gives the first transition. For the second,
the decomposition
$z/(z-\zeta)=1+\zeta/(z-\zeta)$ leaves only the simple
fraction after Cartier, giving the second formula.
Each map preserves the compatibility equation, since
Frobenius and its inverse commute and
$(\zeta b_\zeta)^2=\zeta^2b_{\zeta^2}$.
They are $\Ftwo$-linear maps of the finite profile space.

The constant coefficient is $R(0)$, which is the displayed
sum because the characteristic is two. Frobenius permutes
its summands, so the sum lies in $\Ftwo$.
The profiles reconstruct the entire rational state, and
the two transition equations agree with the two binary
selectors by (L). Induction on the remaining word gives
the final assertion.
\end{proof}

If $S_d$ is the squarefree part of $A_d$, the profile space
is also the space of fractions $U/S_d$ with
$U\in\Ftwo[z]$, $\deg U<\deg S_d$. It therefore has
exactly $2^{\deg S_d}$ states. The initial transient can
be computed in $V_d$ by the numerator matrices and then
converted to this profile by partial fractions. Neither
representation suppresses a pole, a coefficient, or a
cancellation.

\begin{corollary}[Scalar output after repeated Cartier]
\label{cor:scalar-profile}
Let $\beta_\zeta$ denote the initial coefficient of
$(z-\zeta)^{-1}$ in the partial-fraction expansion of a
proper differential over $\Ftwo$ with denominator $A_d$.
For every $\ell\ge a$,
\[
 [z^0\,dz]\,\mathcal C^\ell(R\,dz)
   =\sum_{\zeta\in\mathcal R_d}
              \frac{\beta_\zeta^{\,2^{-\ell}}}{\zeta}.
 \tag{O}
\]
Here $2^{-\ell}$ denotes $\ell$ applications of inverse
Frobenius in $K$.
\end{corollary}
\begin{proof}
A higher pole never becomes simple under repeated
$\mathcal C$: while it survives its new order is
$1+(k-1)/2$, still greater than one. By time $a$ every
such term has vanished. The initial simple residues have
only undergone inverse Frobenius. Apply the output
formula (U).
\end{proof}

Grouping the terms in (O) by Frobenius orbit expresses the
answer as a sum of finite-field traces. If $\zeta$ has
orbit length $h$, compatibility puts its summand in
$\mathbb F_{2^h}$ and the orbit contributes
$\operatorname{Tr}_{\mathbb F_{2^h}/\Ftwo}
(\beta_\zeta^{2^{-\ell}}/\zeta)$.
This makes the scalar cancellation explicit.
A Frobenius-stable subfamily of poles has autonomous
updates under (U), but generally does not determine
the sum over the omitted poles.

\subsection{The actual observation lies beyond the transient}

For the linked values in \eqref{eq:delta}, $C=2^\ell$ with
$\ell\ge B+1$ by \eqref{eq:growth}. Since
$\mu\le d(1+\lfloor\log_2 d\rfloor)\le B(1+\log_2B)$,
we have $a\le B$ for $B\ge16$: use
$\log_2 B\le B/2$ and $B^2\le2^B$ in this range.
Also $B>2u$ gives $d=B-u+3>u$ and
$e-1=u-5<D$, so the actual initial numerator is proper.
Thus the actual selected observation
is beyond the higher-pole transient. This pays an actual parameter
bound; it is not an unspecified promise that enough iterations occur.
Together, \eqref{eq:selection} and (O) give the exact
scalar formula for that linked coefficient using all its
initial simple residues.

The complete state can nevertheless have a long simple-pole cycle.
Pole count alone cannot prove that. For example, $dS/S$ for a
squarefree polynomial $S$ has residue one at every simple pole and
is Cartier-fixed, however many poles $S$ has.


\section{An exact period on the actual linked inputs}

Retain $u,B,C,e,d$ from \eqref{eq:delta}, and write
\[
 B=2^K,\qquad r=K-1,\qquad q=2^r+1=B/2+1,\qquad
 m_0=u-2,\qquad \Delta=m_0(m_0+1)/2.
\]
The integer $q$ need not be prime. All primitive $q$-th roots
in an algebraic closure of $\Ftwo$ will be used.

\begin{lemma}[The actual residue formula]\label{lem:residue}
The form $\omega_e=z^{e-1}A_d(z)^{-1}\,dz$ has a simple nonzero
residue at each primitive $q$-th root $\zeta$, given by
\[
 b_\zeta=\frac{\zeta^e}{A_{d-q}(\zeta)}
                 =\zeta^{e-\Delta}A_{m_0}(\zeta).
\]
\end{lemma}
\begin{proof}
The hierarchy bounds give $d/2<q\le d<2q$ and $m_0<q$.
Thus only $1+z^q$ among the factors of $A_d$ vanishes at
$\zeta$, and its derivative there is $\zeta^{-1}$.
The product $\prod_{i=1}^{q-1}(1+\zeta^i)$ is one: it is
the derivative at one of $X^q-1$, whose value is $q=1$ in
characteristic two. It follows that the residue is the first
displayed expression.

Now $d-q=q-1-m_0$. The last $m_0$ terms in that complete
product equal
$\prod_{j=1}^{m_0}(1+\zeta^{-j})
  =\zeta^{-\Delta}A_{m_0}(\zeta)$.
Taking the complementary product proves the second expression.
None of $1+\zeta^j$ for $1\le j\le m_0$ vanishes.
\end{proof}

\begin{lemma}[The bounds needed for the root count]\label{lem:bounds}
For the actual linked inputs, $r\ge96$, $\Delta<r$, and
\[
 0<\Delta-2e<q,\qquad
 \varphi(q)>\frac{2^r}{r+1}>2^{2r/3}\Delta.
\]
\end{lemma}
\begin{proof}
For $s=u+100$, $F_5(s)\ge F_1(s)\ge2^{2^s}$, so
$K\ge2^{u+100}$ and $r\ge2^{u+100}-1$.
These inequalities imply $r\ge96$ and
$\Delta=(u-2)(u-1)/2<r$.
As $e=u-4=m_0-2$,
$\Delta-2e=(m_0^2-3m_0+8)/2$ is positive and smaller
than $q$ by the same growth bound.

If an odd integer $q$ has $k$ distinct prime factors
$p_1<\cdots<p_k$, then $p_i\ge i+2$, hence
$q/\varphi(q)=\prod p_i/(p_i-1)\le k+1$.
Here $3^k\le q=2^r+1$ implies $k\le r$, so
$\varphi(q)\ge q/(r+1)>2^r/(r+1)$.
Finally $2^{r/3}>r^3$ at $r=96$ and thereafter:
the ratio of the exponential to $r^3$ increases for $r\ge96$.
Since $\Delta<r$ and $r(r+1)<r^3$, this gives
$2^r/(r+1)>2^{2r/3}\Delta$.
\end{proof}

\begin{theorem}[Exact actual residue-tuple period]\label{thm:actual-period}
The tuple $(b_\zeta)_{\zeta\ {\rm primitive}\ q}$ has exact
Frobenius period $2r$. Every eventual period of the complete
rational differential Cartier orbit of $\omega_e$ is divisible
by $2r=2(\log_2 B-1)$.
\end{theorem}
\begin{proof}
The order of two modulo $q=2^r+1$ is $2r$.
It divides $2r$ because $2^r\equiv-1\pmod q$.
Any proper divisor of $2r$ is at most $r$, and for
$s\le r$ the positive integer $2^s-1$ is smaller than $q$.
Thus no smaller order is possible. The primitive roots and their
residues belong to $\mathbb F_{2^{2r}}$, so the tuple period
divides $2r$.

Reciprocity gives
$A_{m_0}(\zeta^{-1})=\zeta^{-\Delta}A_{m_0}(\zeta)$.
If Frobenius power $r$ fixed the tuple,
Lemma~\ref{lem:residue} would give
$\zeta^{\Delta-2e}=1$ at a primitive root, contrary to
$0<\Delta-2e<q$.

Any other proper divisor $s$ of $2r$ is at most $2r/3$:
the integer quotient $(2r)/s$ is then at least three.
Fixedness at all primitive roots would make every such root a zero of
\[
 P_s(X)=A_{m_0}(X)^{2^s}
       +X^{(\Delta-e)(2^s-1)}A_{m_0}(X).
\]
This polynomial has constant coefficient one and degree exactly
$2^s\Delta$, since $e>0$ makes the second summand's degree
strictly smaller. It cannot have all $\varphi(q)$ primitive roots
as zeros, by Lemma~\ref{lem:bounds}. Thus no proper divisor is a
tuple period.

The chosen finite poles were already simple initially.
By Lemma~\ref{lem:local}, the tuple in the $j$-th Cartier state is
$(b_\zeta^{\,2^{-j}})_\zeta$ with poles fixed. If states at
$j$ and $j+p$ agree, these tuples agree. Inverse Frobenius is
bijective and its tuple orbit has period $2r$, so $2r\mid p$.
This applies also after the higher-pole transient.
\end{proof}

The proof is a tuple argument. It does not assert that every individual
residue has degree $2r$, and it does not calculate the selected
constant coefficient. In particular, the growing tuple period does
not prove that the sequence $v\mapsto\delta_4(v)$ is nonperiodic.



\section{Consequences and the remaining scalar question}

Theorem~\ref{thm:transfer} proves the count transfer in
Theorem~\ref{thm:main}. Theorems~\ref{thm:observer}
and~\ref{thm:residue-observer}, together with
Lemma~\ref{lem:observable}, supply its two complete observer
representations. The explicit boundary calculation and
Theorem~\ref{thm:actual-period} supply its linked-input period claim.
Every one of these assertions is about the same specified count or
rational coefficient.

The cutoff-dependent observer is effective for every input. Its size
grows with $d$, and the actual full-state period is not bounded
independently of the entry parameter. Thus attaching an arbitrary
finite automaton to each cutoff is not, by itself, an algorithm for a
quantified statement over all linked inputs.

A smaller observer might preserve only the scalar outputs allowed by
the prescribed input family. All-continuation observability does not
exclude that possibility: its distinguishing words range over every
binary continuation. Likewise, the primitive-root period theorem
does not control cancellation in the sum over all poles. It is not a
proof that $v\mapsto p_d(C-e)\bmod2$ is nonperiodic.

Complete descent counts offer another possible route. The first
retained-prefix block is one summand in a longer count. Observability
of that full sum need not define the graph of its first summand.
The present results therefore provide a precise unchanged-control
arithmetic problem and an exact finite calculation for each input,
without asserting that this isolated bit is a necessary intermediate
for every canonical-comparison algorithm.

Controls chosen to enforce simpler count residues constitute a different
construction. The general transfer theorem can be applied to them when
its congruence premise is proved, but the standard-control profile
theorems here neither transfer automatically to those controls nor
obstruct their comparison algorithms.

\paragraph{Preparation and verification.}
This manuscript was prepared with AI assistance and presents hand
proofs. The finite algebraic checks accompanying the manuscript do not
constitute a formal verification of the ordinal or quantified results.
Historical source dependencies and the relationship to the companion
articles are recorded separately.

\bibliographystyle{plain}
\bibliography{references}
\end{document}
